\chapter{一袋零件}
% 完整版：chapters/01b_neurontypes.tex

\pencilfig{1}{\pencilart{1}}{一袋零件：几乎都是蓝色的，少量红色的，还有一小撮别的。}

假设有人给你一个袋子，里面装着八百六十亿个外表一模一样的零件，请你弄明白这台计算机是怎么用它们搭起来的。你会从哪里下手？工程师不会去端详零件的形状。他会问：这个零件有几个输入、几个输出，输出的是什么。

这种零件叫神经元，它的构造出奇地简单。它有成千上万个输入：来自其他神经元的导线都接到它身上。输出只有一个。但这个输出像树一样分叉，同一个信号的副本会送到成千上万个接收者那里。想象一个人同时听一千个人说话，只回答一句话，而这句话又会被另外一千个人听到。

这是什么样的信号？是短促的“咔嗒”一声。神经元时不时发出一个脉冲——一次电压的突跳，长约千分之一秒。同一个神经元的每一声咔嗒都一样，就像同一面鼓的每一次敲击。所以，神经元想说的一切都不在咔嗒的响度里，而在它\emph{什么时候}咔嗒：更频繁还是更稀疏，立刻还是延迟。

神经元内部做的是能想到的最简单的计算：把传进来的信号加起来，一旦总和越过阈值，自己就咔嗒一声。有些输入把总和往上推，另一些往下拉。

\section*{怎样给这袋零件分类}

输出导线的末端伸到邻居那里，却并不接触它。两者之间留有一道窄缝，信号靠化学方式跨过去：导线释放出一小撮特殊物质——神经递质——邻居再把它接住。幸运的是：几乎每个神经元在自己导线的所有末端释放的都是同一种物质。因此，可以按神经元释放\emph{什么}来给它们分类。

我们用神经元所释放物质的英文名的前四个字母来称呼它的类型。释放谷氨酸（glutamate）的神经元叫\nt{GLUT}。释放$\gamma$-氨基丁酸（GABA）的叫\nt{GABA}。接下来还有\nt{DOPA}（多巴胺）、\nt{SERO}（血清素）、\nt{NORA}（去甲肾上腺素）、\nt{ACET}（乙酰胆碱）等几种。

分好类以后，你会看到一幅极不均衡的图景。每一千个神经元里，大约九百六十个是\nt{GLUT}。它们的信号推动接收者去咔嗒：这是“加号”。还有大约三十六个是\nt{GABA}，它们的信号把接收者推离阈值：这是“减号”。在大脑皮层里，这类神经元更多，占五分之一到四分之一。而其余所有类型加在一起，每十万个神经元里只有几个。沧海一粟。

\section*{电话与电台}

但这一小撮的工作方式截然不同。\nt{GLUT}或\nt{GABA}神经元的导线通向特定的邻居，就像电话线：只有被拨通的人才收到信号。而\nt{DOPA}、\nt{SERO}、\nt{NORA}和\nt{ACET}这些小小的群体则像电台。它们的导线铺满大脑的大片区域，物质常常不是释放到窄缝里，而是释放到周围空间中，同一条消息会同时被几百万个细胞听到。

电话线传的是数据：这是画面上的一条线，那是某个音高的声音。电台播的是全局设置：整体音量调多大，现在值不值得学习，是不是该睡觉了。任何计算机都是这样：存储单元有几十亿个，控制寄存器却只有一小撮。可没有它们，机器照样转不起来。

\section*{含义由接收者决定}

这里有个微妙之处。物质本身并不知道自己是“加号”还是“减号”。分子是一把钥匙。钥匙插进去以后会发生什么，由锁决定——锁是接收者一侧的一种特殊蛋白质。同样是多巴胺，对有些接收者是鼓励，对另一些是抑制，就看它们装的是哪种锁。这就像同一档广播节目，不同的听众听出不同的意思。

\section*{“比预期的好”这个信号}

最著名的广播频道是多巴胺。它传的是什么？来看一个实验。突然给动物一滴果汁，多巴胺神经元会报以一阵密集的咔嗒。之后每次给果汁前都先点亮一盏灯。动物学会了灯意味着果汁，于是这阵爆发改为出现在灯亮时，而果汁本身不再引起爆发。如果灯亮之后没有给果汁，多巴胺神经元在预期的那一刻会\emph{沉默下来}。

\pencilfig{101}{%
% three rows: a spike raster of a DOPA neuron; lamp (amber) and juice (blue drop) above the axis
\foreach \row/\y in {1/0,2/-1.05,3/-2.1}{
  \draw[pen] (0,\y) -- (6.2,\y);
  \foreach \s in {0.3,0.75,1.1,2.15,2.6,3.05,3.45,3.9,4.45,4.95,5.4,5.85}{
    \pgfmathtruncatemacro{\gap}{(\row==3)&&(\s>3.6)&&(\s<4.7)}
    \ifnum\gap=0 \draw[pcGray, line width=0.7pt] (\s,\y) -- (\s,\y+0.3);\fi}}
% row 1: juice without a lamp -> burst at the juice
\foreach \s in {4.0,4.08,4.16,4.24,4.33}{\draw[pcAmber, line width=0.8pt] (\s,0) -- (\s,0.45);}
\draw[dotpen=pcBlue, wash=pcBlue] (4.15,0.72) circle (0.09); \draw[pcBlue, line width=0.6pt] (4.07,0.76) -- (4.15,0.92) -- (4.23,0.76);
% row 2: lamp -> burst at the lamp, nothing at the promised juice
\foreach \y in {-1.05,-2.1}{
  \foreach \s in {1.5,1.58,1.66,1.74,1.83}{\draw[pcAmber, line width=0.8pt] (\s,\y) -- (\s,\y+0.45);}
  \draw[dotpen=pcAmber, wash=pcAmber] (1.65,\y+0.72) circle (0.1);
  \foreach \a in {30,90,150}{\draw[pcAmber, line width=0.5pt] ($(1.65,\y+0.72)+(\a:0.14)$) -- ($(1.65,\y+0.72)+(\a:0.24)$);}}
\draw[dotpen=pcBlue, wash=pcBlue] (4.15,-0.33) circle (0.09); \draw[pcBlue, line width=0.6pt] (4.07,-0.29) -- (4.15,-0.13) -- (4.23,-0.29);
% row 3: lamp, no juice -> a gap where the juice was expected
\draw[pcBlue, line width=0.6pt, densely dotted] (4.15,-1.38) circle (0.09); \draw[pcBlue, line width=0.6pt, densely dotted] (4.07,-1.34) -- (4.15,-1.18) -- (4.23,-1.34);
\draw[penlite=pcRed] (3.65,-2.22) -- (3.65,-2.28) -- (4.65,-2.28) -- (4.65,-2.22);
\node[lbl, text=pcRed, anchor=north] at (4.15,-2.3) {低谷};
\node[lbl, anchor=east, align=right] at (-0.1,0.15) {果汁，\\没有灯};
\node[lbl, anchor=east, align=right] at (-0.1,-0.9) {灯亮，\\然后果汁};
\node[lbl, anchor=east, align=right] at (-0.1,-1.95) {灯亮，\\没有果汁};
}{对意外的果汁爆发；对预告果汁的灯爆发；在预告的果汁没有到来的那一刻出现低谷。}

有一条规则能解释全部三种情况：多巴胺报告的不是“好”，而是“比我预期的好”。意外的果汁是惊喜。预告过的果汁没有任何意外。预告了却没给，是负面的意外。那些从自身错误中学习的程序，用的正是这个量。我们后面还会回到它。

\section*{带走什么}

大脑是一张巨大的网络，由外表相同的零件组成。它们几乎全是数据电话网里的“加号”和“减号”。只有极少数是电台，一次调节整张网络。下一章我们来看看，只用“加号”能搭出什么。

\fullversion{1}
